\chapter*{Anexo B}
\addcontentsline{toc}{chapter}{Anexo B: Procesamiento de cargas del DLC 1.1}

\setcounter{table}{0}
\renewcommand{\thetable}{B.\arabic{table}}
\setcounter{figure}{0}
\renewcommand{\thefigure}{B.\arabic{figure}}

\section*{Procesamiento de cargas del DLC 1.1: extrapolación al período de retorno de 50 años}
\label{anexo:dlc11}

Este anexo reproduce el código que transforma las series temporales entregadas por QBlade en las cargas de diseño del DLC 1.1, siguiendo el procedimiento del Anexo G de IEC 61400-1~\parencite{IEC61400-1} descrito en la Sección~\ref{sec:anexo_f}. Se omiten las rutinas de exportación a planilla y de generación de figuras, por no formar parte del cálculo. El procesamiento parte de $72$ realizaciones, correspondientes a $12$ velocidades de viento y $6$ semillas de turbulencia cada una.

\subsection*{Geometría y parámetros de entrada}

La pala se discretiza en los mismos $21$ paneles de reparto coseno con que se resuelve la aerodinámica, de modo que la carga se lee donde el método calcula efectivamente la circulación (Sección~\ref{sec:discretizacion}).

\begin{pycode}
# -*- coding: utf-8 -*-
"""
DLC 1.1 - Extrapolación de cargas extremas según el Anexo G de IEC 61400-1.
Entrada : 72 series temporales de fuerzas por estación (12 V x 6 semillas).
Salida  : carga característica F_k a 50 años y carga de diseño F_d de cada
          una de las trece magnitudes, y el perfil espacial de diseño.
"""
import csv, glob, json, math, os, re
from collections import defaultdict
import numpy as np

# --- Geometría de la pala y discretización ------------------------------
# Los bordes de panel siguen el reparto coseno, que concentra estaciones
# hacia los extremos, donde los gradientes de circulación son mayores. Los
# anchos suman por construcción la envergadura exacta, de modo que la
# integración espacial no introduce error de cierre.
ENVERGADURA = 23.0
N_PAN   = 21
BORDES  = 0.5 * ENVERGADURA * (1.0 - np.cos(np.linspace(0.0, np.pi, N_PAN + 1)))
Z_PAN   = 0.5 * (BORDES[:-1] + BORDES[1:])   # centro de panel [m]
DELTA_S = np.diff(BORDES)                    # ancho de panel [m]

# Secciones de apoyo: las dos conexiones con los brazos, que son la "raíz"
# de la pala en el sentido de la Sección 3.1 de IEC 61400-5.
Z_APOYO_INF = 5.0
Z_APOYO_SUP = 18.0

# Direcciones de la envolvente. IEC 61400-5 6.1.4.3 da por suficiente un
# conjunto repartido uniformemente en al menos doce direcciones.
N_DIRECCIONES = 12

# --- Parámetros normativos ----------------------------------------------
T_DISCARD = 50.0                       # tramo inicial descartado [s]. La
                                       # norma exige eliminar al menos los
                                       # 5 s iniciales, o mas si es
                                       # necesario (Seccion 7.5).
T_OBS     = 600.0                      # ventana útil de cada realización [s]
T_RETURN  = 50 * 365.25 * 24 * 3600.0  # período de retorno de 50 años [s]

# Probabilidad de excedencia objetivo. La Sección 7.6.2.2 de la norma
# declara ese mismo valor, 3,8e-7, como el máximo admisible.
PE_TARGET = T_OBS / T_RETURN

# Factor parcial de carga. La norma admite reducirlo de 1,35 a 1,25 para
# el DLC 1.1 cuando la carga se obtiene por extrapolación estadística
# según el Anexo G (Sección 7.6.2.2, nota a de la Tabla 3).
GAMMA_F = 1.25
\end{pycode}

\subsection*{Efectos de carga que se extrapolan}

QBlade entrega, en cada paso de tiempo, la fuerza por unidad de longitud en las $21$ estaciones y en las dos direcciones del sistema local de la pala contenidas en el plano del perfil. De ese vector se derivan los efectos de carga sobre los que opera la extrapolación, que son los momentos flectores en las dos secciones de apoyo.

\begin{pycode}
def fuerzas_puntuales(fn_lin, ft_lin):
    """Pasa de fuerza por unidad de longitud [N/m] a fuerza de panel [N]."""
    return fn_lin * DELTA_S, ft_lin * DELTA_S


def momentos_en_apoyo(fn_lin, ft_lin, z_c, lado):
    """
    Momento flector en una sección de apoyo, producido por el voladizo que
    esa sección sostiene.

    Cortando en el apoyo y sumando solo el voladizo libre, el resultado
    queda estáticamente determinado y no depende de las condiciones de
    contorno. El vano central entre apoyos no lo está, y por eso se
    resuelve en el modelo de elementos finitos y no aquí.

    Devuelve (m_x, m_y): m_x lo genera la componente tangencial y flecta
    en torno al eje local X; m_y lo genera la normal y flecta en torno a Y.
    """
    fn, ft = fuerzas_puntuales(fn_lin, ft_lin)
    if lado == 'inf':
        sel   = Z_PAN < z_c
        brazo = z_c - Z_PAN[sel]
    else:
        sel   = Z_PAN > z_c
        brazo = Z_PAN[sel] - z_c
    m_x = (ft[..., sel] * brazo).sum(axis=-1)
    m_y = (fn[..., sel] * brazo).sum(axis=-1)
    return m_x, m_y


def envolvente_direccional(m_x, m_y, n=N_DIRECCIONES):
    """
    Máximo del momento proyectado sobre n direcciones repartidas
    uniformemente en la sección (IEC 61400-5 6.1.4.3, enfoque avanzado).

    Al recorrer las n direcciones el círculo completo, la proyección cubre
    los dos sentidos de cada componente y no hace falta tomar valor
    absoluto. Por construcción el resultado nunca supera el módulo del
    vector, y solo lo iguala cuando este apunta a una de las direcciones.
    """
    th   = np.linspace(0.0, 2.0 * math.pi, n, endpoint=False)
    proy = (np.outer(np.atleast_1d(m_x), np.cos(th))
            + np.outer(np.atleast_1d(m_y), np.sin(th)))
    return proy.max(axis=1)


def f_total(fn_lin, ft_lin):
    """Fuerza total integrada sobre la envergadura, en norma euclídea."""
    fn, ft = fuerzas_puntuales(fn_lin, ft_lin)
    return np.hypot(fn, ft).sum(axis=-1)


def efectos_de_carga(fn_lin, ft_lin):
    """Las trece magnitudes que se extrapolan, para una serie completa."""
    mx_i, my_i = momentos_en_apoyo(fn_lin, ft_lin, Z_APOYO_INF, 'inf')
    mx_s, my_s = momentos_en_apoyo(fn_lin, ft_lin, Z_APOYO_SUP, 'sup')
    return {
        # Componentes con signo, en sus dos sentidos: el extremo de diseño
        # de una magnitud predominantemente negativa está en su mínimo, y
        # la norma pide los valores extremos, que son ambos.
        'M_x_apoyo_inf_pos': mx_i, 'M_x_apoyo_inf_neg': -mx_i,
        'M_y_apoyo_inf_pos': my_i, 'M_y_apoyo_inf_neg': -my_i,
        'M_x_apoyo_sup_pos': mx_s, 'M_x_apoyo_sup_neg': -mx_s,
        'M_y_apoyo_sup_pos': my_s, 'M_y_apoyo_sup_neg': -my_s,
        # No negativas por construcción.
        'M_res_apoyo_inf': np.hypot(mx_i, my_i),
        'M_env_apoyo_inf': envolvente_direccional(mx_i, my_i),
        'M_res_apoyo_sup': np.hypot(mx_s, my_s),
        'M_env_apoyo_sup': envolvente_direccional(mx_s, my_s),
        'F_tot': f_total(fn_lin, ft_lin),
    }
\end{pycode}

\subsection*{Ajuste de la distribución de extremos por velocidad}

De cada realización se extrae el máximo de cada magnitud sobre la ventana útil. Para cada velocidad de viento se dispone así de seis máximos, uno por semilla, a los que el Anexo G ajusta una distribución de Gumbel por el método de los momentos, según las Ecs.~\ref{eq:gumbel_beta} y~\ref{eq:gumbel_mu}.

\begin{pycode}
GAMMA_EULER = 0.5772156649   # constante de Euler-Mascheroni

def ajuste_gumbel(maximos):
    """
    Ajusta una distribución de Gumbel por el método de los momentos, tal
    como prescribe la Sección G.3.2 del Anexo G:

        beta = sqrt(6)/pi * s          mu = x_medio - gamma_E * beta

    donde s es la desviación estándar muestral de los máximos y x_medio su
    promedio.
    """
    x    = np.asarray(maximos, dtype=float)
    beta = np.sqrt(6.0) / np.pi * np.std(x, ddof=1)
    mu   = np.mean(x) - GAMMA_EULER * beta
    return mu, beta


def excedencia_gumbel(F, mu, beta):
    """
    Probabilidad de que la magnitud supere el valor F en una realización
    de 600 s a una velocidad de viento dada:

        P(S > F | V) = 1 - exp( -exp( -(F - mu) / beta ) )
    """
    return 1.0 - np.exp(-np.exp(-(F - mu) / beta))
\end{pycode}

\subsection*{Agregación y carga característica}

La probabilidad de excedencia de largo plazo pondera la probabilidad condicional de cada velocidad por el tiempo que el viento permanece en ella. Esa ponderación emplea la frecuencia empírica de la serie del emplazamiento y no la Weibull ajustada, por las razones expuestas en la Sección~\ref{sec:anexo_f}.

\begin{pycode}
# Frecuencia empírica por banda de 1 m/s, contada sobre los 38 años de la
# serie extrapolada a la altura del centro del rotor. La Sección G.3.4
# admite ponderar por la Weibull "o la que resulte apropiada"; aquí la
# apropiada es la empírica, porque la Weibull subestima la banda alta y
# ese sesgo va en contra del margen.
with open('distribucion_viento.json') as fh:
    PESO = {float(k): v for k, v in json.load(fh)['bandas'].items()}

V_IN, V_OUT, DV = 2.0, 25.0, 1.0


def malla_de_integracion(v_sim, mus, betas):
    """
    Bandas sobre las que se integra, y los parámetros de Gumbel
    interpolados en ellas.

    La malla cubre solo el rango operacional. Por debajo de V_in el rotor
    no produce y la interpolación le asignaría la Gumbel de V_in, que es
    una distribución de cargas que ahí no existe; por encima de V_out la
    turbina está detenida y sus cargas son las de los DLC 6.x.
    """
    centros = np.arange(V_IN + DV / 2.0, V_OUT, DV)
    return (centros,
            np.interp(centros, v_sim, mus),
            np.interp(centros, v_sim, betas))


def excedencia_largo_plazo(F, centros, mus, betas):
    """
    Probabilidad de excedencia de largo plazo (Anexo G, Ec. G.2):

        P_e(F) = suma_V  P(S > F | V) * p(V)

    Combina la distribución de extremos condicional a cada velocidad con
    la frecuencia con que esa velocidad ocurre en el emplazamiento.
    """
    return float(sum(excedencia_gumbel(F, mu, beta) * PESO[v]
                     for v, mu, beta in zip(centros, mus, betas)))


def carga_caracteristica(centros, mus, betas):
    """
    F_k por bisección. P_e es monótona decreciente en F, de modo que la
    bisección converge sin ambigüedad al único valor que satisface
    P_e(F_k) = PE_TARGET.
    """
    lo, hi = 0.0, max(mus) + 60.0 * max(betas)
    for _ in range(300):
        med = 0.5 * (lo + hi)
        if excedencia_largo_plazo(med, centros, mus, betas) > PE_TARGET:
            lo = med
        else:
            hi = med
    return 0.5 * (lo + hi)
\end{pycode}

\subsection*{Recorrido de las realizaciones y cargas de diseño}

El procedimiento anterior se aplica de forma independiente a cada una de las trece magnitudes. En el mismo recorrido se localiza el instante gobernante, que es aquel en que la magnitud elegida para escalar el perfil alcanza su máximo sobre el conjunto de las 72 realizaciones.

\begin{pycode}
MAGNITUD_GOBERNANTE = 'M_env_apoyo_sup'   # sección más solicitada

maximos = defaultdict(lambda: defaultdict(list))
mejor   = {'valor': -np.inf}

for ruta in sorted(glob.glob('sim_DLC11_*_forces.csv')):
    # Nombre del archivo: sim_DLC11_068_V24_S2_forces.csv -> V = 24, semilla 2
    m = re.search(r'sim_DLC11_(\d+)_V(\d+)_S(\d+)_forces\.csv', ruta)
    v, semilla = int(m.group(2)), int(m.group(3))

    t, fn, ft = leer_csv(ruta)          # fn, ft en [n_pasos, N_PAN], N/m
    util = t >= T_DISCARD               # descarte del transitorio inicial
    ef = efectos_de_carga(fn[util], ft[util])

    for nombre, serie in ef.items():
        maximos[nombre][v].append(float(serie.max()))

    # Instante gobernante, para el perfil espacial de diseño.
    serie = ef[MAGNITUD_GOBERNANTE]
    j = int(np.argmax(serie))
    if serie[j] > mejor['valor']:
        mejor = {'valor': float(serie[j]), 't': float(t[util][j]),
                 'V': v, 'semilla': semilla,
                 'fn': fn[util][j].copy(), 'ft': ft[util][j].copy()}

resultados = {}
for nombre, por_v in maximos.items():
    vs    = sorted(por_v)
    pares = [ajuste_gumbel(por_v[v]) for v in vs]
    centros, mu_i, beta_i = malla_de_integracion(
        np.array(vs, float),
        np.array([p[0] for p in pares]),
        np.array([p[1] for p in pares]))
    fk = carga_caracteristica(centros, mu_i, beta_i)
    resultados[nombre] = {'F_k': fk, 'F_d': GAMMA_F * fk}
\end{pycode}

\subsection*{Del escalar a las fuerzas por estación}

El Anexo G entrega un valor escalar por magnitud, que debe convertirse en una distribución de fuerzas a lo largo de la envergadura para cargar el modelo estructural. Para preservar la correlación espacial entre estaciones, el perfil se toma del instante gobernante y se escala de manera uniforme, conforme a la Sección 7.5 de la norma. Las cargas que actúan simultáneamente con el extremo son las que el Anexo I denomina cargas compañeras.

\begin{pycode}
# Valor de la magnitud gobernante en el instante gobernante, y factor
# único que lo lleva hasta la carga de diseño conservando intacta la forma
# del perfil a lo largo de la envergadura.
escala = resultados[MAGNITUD_GOBERNANTE]['F_d'] / mejor['valor']

# Fuerzas de diseño por estación [N]. Se conserva el signo original, de
# modo que la dirección de aplicación en el modelo estructural reproduce
# la del instante gobernante: normal positiva hacia el eje del rotor y
# tangencial positiva en el sentido del movimiento de la pala.
Fd_n = mejor['fn'] * DELTA_S * escala
Fd_t = mejor['ft'] * DELTA_S * escala
\end{pycode}

El mismo cálculo se repite con \texttt{MAGNITUD\_GOBERNANTE} apuntando al apoyo inferior, y con el extremo de $M_y$ en su sentido contrario, que selecciona el instante en que la carga normal apunta hacia afuera y se suma a la centrífuga. Las cuatro distribuciones resultantes se recogen en las tablas del Anexo~A.
